\documentclass[a4paper, 11pt]{article}

% Packages
\usepackage{graphicx}
\usepackage{amsmath}
\usepackage{authblk}
\usepackage{appendix}
\usepackage{listings}
\usepackage{subcaption}
\usepackage{siunitx}
\usepackage[T1]{fontenc} 
\usepackage[a4paper, margin=1in]{geometry}
\usepackage[dvipsnames]{xcolor}

% Hyperlink configuration
\usepackage[hidelinks]{hyperref}
\usepackage[nameinlink, capitalise]{cleveref}

% Code listing configuration
\lstset{ 
    language=Python,               
    keepspaces=true,
    columns=fullflexible,
    backgroundcolor=\color{white},  
    basicstyle=\ttfamily, 
    showstringspaces=false,
    keywordstyle=\color{blue},      
    commentstyle=\color{ForestGreen},     
    stringstyle=\color{red},        
    breaklines=true,                
    captionpos=b,                   
    numbers=left,                   
    numberstyle=\tiny\color{gray}   
}

% Bibliography setup
\usepackage[backend=biber,style=apa]{biblatex}
\addbibresource{references.bib}

\begin{document}

% ==========================================
%         PROJECT PLAN TITLE PAGE
% ==========================================
\begin{titlepage}
    \centering
    \vspace*{1cm}
    
    {\Large\bfseries Project Plan \par}
    \vspace{0.5cm}
    {\large Imperial College London \par}
    {\large Department of Earth Science and Engineering \par} 
    {\large MSc Environmental Data Science and Machine Learning \par} 
    
    \vspace{2.5cm}
    
    {\LARGE\bfseries High-Resolution Spatial Downscaling of Wildfire Data \par}
    
    \vspace{3cm}
    
    {\large\bfseries Author:} \\
    {\large Tsuyu Huang} \\
    {\small \href{mailto:tsuyu.huang25@imperial.ac.uk}{tsuyu.huang25@imperial.ac.uk} \par} 
    
    \vspace{1.5cm}
    
    {\large\bfseries Supervisors:} \\
    {\large Jorge Veiras (\href{mailto:jorge@artiocarbon.com}{jorge@artiocarbon.com})} \\
    {\large Prof. Anna Korre (\href{mailto:a.korre@imperial.ac.uk}{a.korre@imperial.ac.uk}) \par}

    \vspace{1.5cm}
    
    {\large\bfseries Company:} \\
    {\large Artio} \\
    {\small \href{https://www.artiocarbon.com/}{www.artiocarbon.com} \par}
    
    \vfill
    
    {\large \today \par}
\end{titlepage}

\nopagebreak

% ==========================================
%          ABSTRACT & CONTENTS
% ==========================================
\begin{abstract}
Wildfires are destructive and a rapidly evolving hazard. With climate change, they become increasingly severe and frequent. This increases the need for a higher-resolution wildfire data for better monitoring and risk assessment. Geostationary satellites (GEO) have high-temporal resolution, but lacks the resolution needed for a detailed wildfire risk assessment. While low orbit satellites (LEO) provide better spatial resolution, but low temporal resolution to monitor rapid wildfire spread. Existing generative downscaling models have demonstrated their effectiveness in producing sharp, realistic high-resolution remote sensing data. But many still suffers from artefacts, bluriness, and hallucinations. This project aims to achieve a high-resolution spatial downscaling of wildfire data using Conditional Diffusion Models (DDPMs). The model integrates additional physical covariates, such as topography and landcover to mitigate physically implausible outputs. The project compares direct (2km to 30m) and cascaded (2km to 375m to 30m) downscaling pipelines and their effectiveness in producing high-fidelity, high-resolution wildfire data. The project compares basline vanilla DDPM with 3 other DDPMs with different noise prediction backbone, which is based on novel downscaling architectures: Cross Resolution Distance Aware (CredDSR), Terrain-Adaptive U-Net Transformer (TAUT), and .
\end{abstract}

\tableofcontents
\clearpage

% ==========================================
%             MAIN CONTENT
% ==========================================
\section{Introduction}
In recent years, due to rising popularity in Net-Zero initiatives, carbon investors are increasingly funding carbon reduction and removal projects, such as reforestation (\cite{InvestmentTrendsOutcomes}). However, these carbon projects are at risk because wildfires can destroy and release the carbon back into the atmosphere. With climate change, wildfires are becoming increasingly severe and frequent (\cite{jonesClimateChangeIncreases,}). Which increases the need for better wildfire monitoring and prediction for risk assessment. This means high-spatial and high-temporal resolution wildfire data is highly desired. However, high-temporal resolution data obtained from geostationary satellites (GEO) have limited spatial resolution, which makes them difficult to use for accurate risk assessment. While low orbit satellites (LEO) offers high-spatial resolution, but low-temporal resolution.

To bridge this gap, existing methods uses Deep Learning (DL) models for downscaling due to their ability in recognising non-linear spatiotemporal patterns. However, many existing DL downscaling models are often deterministic, and in the context of wildfire risk assessment, probabilistic models are more desirable as they capture the uncertainty and variability in wildfire behavior (\cite{yuProbabilisticApproachWildfire2026}). This project aims to downscale coarse wildfire data using conditional diffusion models (DDPM) due to their probabilistic nature. To mitigate the model from producing physically implausible outputs, the model integrates additional physical information, such as topography and landcover to help guide the downscaling process. The project aims to investigate whether cascaded downscaling (2km to 375m to 30m) produces more accurate and better results compared to direct downscaling (2km to 30m).

\section{Literature Review}

\section{Problem Description and Objectives}

\section{Progress to Date}

\section{Future Plans}

\printbibliography

\end{document}